% !TITLE 狱中题壁
% !AUTHOR 谭嗣同
% !DYNASTY 近代
% !GENRE 七言绝句
% !ORDER 2

\begin{poem}
望门投止思张俭\textsuperscript{1}，忍死须臾待杜根\textsuperscript{2}。
我自横刀向天笑，去留肝胆两昆仑\textsuperscript{3}。
\end{poem}

\begin{annotations}
\notetext{1}{张俭：东汉人，因弹劾宦官被追捕，逃亡途中人们敬重他的品行，纷纷冒死收留他。“望门投止”即看见人家就去投宿（形容逃亡之狼狈）。}
\notetext{2}{杜根：东汉人，因上书触怒邓太后，被装进口袋摔打，假死三日，后来复苏逃匿，做了十五年酒保。忍死须臾——忍住死亡的痛苦，等待片刻（等待复苏的机会）。}
\notetext{3}{肝胆两昆仑：肝和胆（比喻忠心赤胆）像两座昆仑山一样巍峨。谭嗣同赴死前写此诗——去的（赴死）和留的（活着的同志）都如昆仑般顶天立地。}
\end{annotations}

\begin{appreciation}
这是谭嗣同被关在刑部大牢里，题在墙壁上的绝命诗。戊戌变法失败后，他本可以逃走，日本使馆也答应保护他。他拒绝了，说：各国变法没有不流血的，中国还没有人为变法流过血——要有流血的人，就从我谭嗣同开始。几天后，他在菜市口就义，年仅三十三岁。所以这首诗不是写出来的，是一步一步走出来的。

“望门投止思张俭”。张俭是东汉人，因弹劾宦官被追捕，逃亡路上见门就投宿，沿途百姓敬重他的品行，冒死收留他。谭嗣同想起张俭，想的是正在逃亡的战友康有为、梁启超——他们会不会也像张俭一样，一路有人收留？“忍死须臾待杜根”。杜根也是东汉人，因上书触怒太后，被装进口袋摔打，装死逃过一劫，藏匿十五年，终于等到出头之日。谭嗣同是相信的：活着的人，会像杜根一样忍死待时，等来改天换地的那一天。

前两句，他在想别人的活；第三句，轮到自己的死：“我自横刀向天笑”。不逃了。横刀就刑，仰天大笑——这是把生死彻底看透之后才有的笑。最后一联“去留肝胆两昆仑”：去的（他自己）与留的（活着的同志），都是一腔肝胆，都像昆仑山一样巍峨。他把自己和活着的人摆在同一架天平上：留下的人替他接着活，他替留下的人先死。

这个死法，其实是书里读过的那句话的真人版。屈原在《离骚》里说“亦余心之所善兮，虽九死其犹未悔”——只要是我认定的，死九次也不后悔。两千多年来，说过这句话的人很多，把它一句一句活成真的，谭嗣同算一个。中国的诗人不只会写春花秋月，也会写怎么死：屈原沉了江，文天祥过了零丁洋，谭嗣同横刀大笑。一代一代，总有人把“死”的那一份先替大家用掉，剩下的人，才能安心地活。

你当然不用做谭嗣同做的事。把书念好，把日子过好，少熬夜，按时吃饭——那些替你死了的人，要的不过就是这些。
\end{appreciation}
